\section{Meta review of surrogate modeling and multi-energy systems}
\label{sec:related-reviews}

TableÂ~\ref{tab:T7-related-reviews} compares prior review studies
identified through our literature selection process.

\begin{table*}[!t]
\centering
\scriptsize
\begin{tabularx}{\textwidth}{@{}@ L{0.015\textwidth} C{0.025\textwidth} L{0.050\textwidth} L{0.150\textwidth} L{0.155\textwidth} L{0.25\textwidth} L{0.30\textwidth} @ Ref. & Year &@{}}
\toprule
Reviewed material & Method scope & Application scope & Main outcome &
Research gaps identified\
\\
\cite{Bhosekar2018} & 2018 & ref. list: 126 & Surrogate modelling,
feasibility analysis, derivative-free optimization; RBF and kriging &
Generic engineering and chemical-engineering optimization & Reviews
surrogate models for modelling, feasibility analysis and optimization;
emphasizes problem-dependent surrogate choice. & Need for better
model-selection guidance and more reliable surrogate use in constrained
expensive optimization.\
\cite{DiazManriquez2016} & 2016 & ref. list: 64 & Surrogate-assisted MOEAs;
fitness replacement; evolution control & Generic expensive
multi-objective optimization & Classifies surrogate-assisted MOEAs by
how the surrogate interacts with the evolutionary search. & Need to
control surrogate error and avoid false Pareto fronts when true
evaluations are replaced.\
\cite{FernandezGodino2023} & 2023 & ref. list: 116 & Multi-fidelity
surrogate modelling and uncertainty quantification & Generic
engineering, simulation and optimization & Reviews strategies for
combining low- and high-fidelity models to reduce computational cost. &
Need for robust fidelity management, uncertainty propagation and
validation across heterogeneous model fidelities.\
\cite{Starke2025214} & 2025 & 126 selected studies & ML-based
metamodelling and surrogate modelling & Energy systems & Reviews model
applications, ML techniques and error indicators across energy-system
metamodels. & Need for standardized error analysis and comparison of
accuracy together with computational efficiency.\
\\
\cite{Westermann2019} & 2019 & 57 studies & GP, ANN, SVM, RBF, MARS, PCE, RF
and related surrogates & Sustainable building design and building
performance simulation & Establishes surrogate modelling as a practical
emulator for design feedback, sensitivity analysis, uncertainty analysis
and optimization. & Need for careful validation, sampling design and
practical integration into building-design workflows.\
\cite{Elwy2024} & 2024 & 72 papers & ML-based surrogates, sampling
strategies and building-performance optimization & Sustainable building
design optimization & Reviews surrogate-assisted building performance
optimization and identifies dominant variables, objectives, sampling
methods and workflows. & Need for adaptive sampling, better
generalization across climates and more standardized surrogate-based
building optimization.\
\cite{Zhang2026_building} & 2026 & 114 studies & AI-driven surrogates,
ML/DL, MOO algorithms and MCDM & Parametric optimization in sustainable
building design & Proposes a six-step AI-surrogate MOO framework from
dataset construction to optimization and decision-making. & Limited deep
learning use, weak interpretability, insufficient uncertainty treatment
and inadequate block-scale decarbonization analysis.\
\cite{Khan2024_DT} & 2024 & Review / framework & AI-driven surrogate modelling,
uncertainty quantification and digital twins & Hybrid and sustainable
energy systems & Positions surrogate modelling as a computational
accelerator for hybrid energy-system simulation and digital-twin
applications. & Lack of large-scale HES datasets, limited experimental
validation and need for uncertainty-aware surrogate/digital-twin
frameworks.\
\cite{Elsheikh2019622} & 2019 & ref. list: 166 & Artificial neural networks;
MLP, WNN, RBF and Elman networks & Solar energy devices & Reviews ANN
applications for predicting and optimizing solar-device performance. &
Mainly device-level ANN prediction; limited treatment of system-level
optimization and surrogate integration.\
\cite{Etghani2025} & 2025 & ref. list: 170 & ANN, Taguchi methods, GRA and
TOPSIS & Engines, thermal storage, solar/wind systems, heat exchangers
and HRES & Synthesizes ANN--Taguchi approaches for prediction and robust
parameter optimization. & Need for broader ANN--Taguchi frameworks,
uncertainty treatment and validation in complex energy systems.\
\\
\cite{Ruan2021221} & 2021 & 62 selected studies & ML-assisted optimization;
surrogate models and hybrid models & Power-system optimization,
including OPF, UC and economic dispatch & Reviews how ML can assist
power-system optimization and distinguishes prediction, surrogate and
hybrid optimization roles. & Need for deeper ML--optimization
integration, reliability guarantees and scalable deployment in volatile
power systems.\
\cite{Khaloie2025} & 2025 & ref. list: 180 & E2E OPF learning and L2O;
supervised, unsupervised, RL and graph learning & Transmission-level
optimal power flow & Provides a detailed taxonomy and evaluation
framework for learning-based OPF. & Challenges include feasibility,
generalization, topology changes, scalability and reliable
benchmarking.\
\cite{Mohammadi2024} & 2024 & ref. list: 106 & Analytical and learned
surrogate models for OPF & Optimal power flow at transmission and
distribution levels & Classifies OPF surrogates into analytical and
learned surrogates. & Need for stronger constraint handling, topology
generalization and hybrid structures preserving hard OPF constraints.\
\cite{Tan2026} & 2026 & ref. list: 152 & Gaussian processes and uncertainty
quantification & Power-system forecasting, modelling, risk assessment,
optimization and control & Reviews GP fundamentals and power-system
applications with emphasis on uncertainty-aware prediction. & Key
challenges are GP scalability, robustness and integration with modern
data-driven methods.\
\cite{Lim2025} & 2025 & ref. list: 284 & Deep learning; learning to predict,
surrogate and optimize & Power-system decision-making & Classifies
neural networks by functional role in power-system decisions. & Need for
modular/hybrid architectures, domain-specific foundation models and
clearer role-specific evaluation.\
\\
\cite{Manco2024} & 2024 & 6 case studies & LP, MILP, NLP, MINLP;
uncertainty, flexibility, storage and thermal networks & Multi-energy
systems, energy hubs and polygeneration & Provides a
formulation-oriented guide showing how modelling choices affect problem
class, solver selection and complexity. & Need for explicit mathematical
formulations, better thermal-network/storage modelling and scalable MES
optimization.\
\cite{Klemm2021} & 2021 & 145 tools & MILP, LP, dynamic modelling and
bottom-up/hybrid approaches & MES optimization in mixed-use urban
districts & Identifies model requirements and shows that only few tools
are suitable for mixed-use district MES optimization. & Many tools lack
optimization, sector coverage, suitable spatial scale, temporal
resolution or non-financial criteria.\
\cite{Fattahi2020} & 2020 & 19 national ESMs & Integrated ESMs, MCA,
hard/soft linking, optimization and simulation suites & National
low-carbon energy-system models & Identifies seven low-carbon modelling
challenges and maps them to required model capabilities. & Current ESMs
inadequately capture flexibility, decentralization, emerging
technologies, macroeconomic feedbacks and social behaviour.\
\cite{Mylonopoulos202332697} & 2023 & ref. list: 88 & Physics-based and
data-driven modelling; design, EMS/control and combined optimization &
Ship power and propulsion energy systems & Reviews modelling and
optimization of ship energy systems and identifies surrogate-assisted
optimization as promising. & Limited data-driven and surrogate-assisted
optimization literature for integrated ship energy systems.\
\cite{agha_kassab_comprehensive_2024} & 2024 & ref. list: 144 & Heuristic,
mathematical and hybrid optimization; sizing and EMS & Microgrid
planning with PV and renewables & Reviews microgrid sizing and
energy-management strategies and stresses integrated sizing--EMS
planning. & Need for integrated sizing--EMS methods, more MOO, better
degradation/load models and EV/smart-grid integration.\
\cite{nallolla_multi-objective_2023} & 2023 & ref. list: 157 &
Multi-objective optimization algorithms for AC/DC and hybrid microgrids
& Hybrid AC/DC microgrids with renewable energy & Reviews MOO algorithms
and objective functions for hybrid AC/DC microgrid planning and
operation. & Need for better algorithm selection, renewable
intermittency handling and practical validation of MOO-based designs.\
\cite{salgueiro_multi-objective_2019} & 2019 & n.r. & Multi-objective
metaheuristics; NSGA, MOPSO, MOEA/D & Microgrid planning and management
& Surveys technical challenges in applying multi-objective
metaheuristics to microgrid planning and operation. & Need for robust
algorithm design, better performance assessment and stronger treatment
of uncertainty.\
\cite{malla_sg_optimization_2024} & 2024 & n.r. & MILP, MOO and Power-to-X
optimization methods & MES with Power-to-X, hydrogen and district-scale
applications & Reviews optimization approaches for sector-coupled MES
with Power-to-X technologies. & Need for stronger integration of
Power-to-X flexibility, storage and sector-coupled operational
constraints.\
\cite{Zhou2024__2} & 2024 & ref. list: 147 & Sizing methods, MOO, Monte
Carlo simulation and AI-assisted sizing & Centralized and distributed
renewable-storage systems & Compares renewable-storage sizing approaches
and identifies AI-assisted sizing as emerging. & Need for
climate-adaptive sizing, predictive energy management and
lifecycle-oriented storage sizing.\
\cite{Li2025_hydrogen} & 2025 & ref. list: 103 & Mathematical models,
simulation tools, intelligent optimization and physics-assisted AI &
Hydrogen-integrated hybrid renewable energy systems & Reviews
configurations, models, tools and AI-driven optimization for
hydrogen-based HRES. & Need for better data availability,
interpretability, stochastic/physics-assisted AI and broader
social-environmental factors.\
\cite{vahidinasab_overview_2020} & 2020 & ref. list: 184 & Planning models,
single-/multi-objective optimization, uncertainty handling and
decomposition & Distribution network expansion planning with DERs &
Reviews DEP objectives, constraints, horizons, variables, DER expansion
and uncertainty treatment. & Need for integrated planning linking DER
expansion, uncertainty, MOO and district-energy interactions.\
\cite{Conti2026} & 2026 & ref. list: 158 & Forecasting and optimization;
deterministic, stochastic, metaheuristic and hybrid methods &
Large-scale PV integration in smart distribution networks & Reviews
optimization strategies for PV integration and links forecasting, ESS
coordination and MOO. & Need for scalable optimization under
uncertainty, better grid constraints and deployment-oriented
validation.\
\cite{batista_optimizing_2023} & 2023 & 374 records; 38% selected &
Bibliometric analysis, systematic review, Portfolio Theory and PSO &
HRES-powered reverse-osmosis desalination & Maps optimization approaches
for HRES-powered RO plants and proposes PT--PSO for energy management. &
Need for integrated HRES--RO optimization and better treatment of
storage, reliability and water-energy trade-offs.\
\\
\cite{ChenRenZhou2023} & 2023 & Review & Prior, posterior and
interactive MOO methods; preference-based classification & Long-term
energy-system models & Classifies objectives into economic,
environmental, social and energy-security dimensions. & Need for more
social/security objectives, interactive methods and comparisons by
accuracy and efficiency.\
\cite{Sahoo2025} & 2025 & ref. list: 41 & MCDM methods including AHP,
TOPSIS, VIKOR, fuzzy and hybrid models & Energy management, renewable
energy, grid management and policy planning & Reviews MCDM applications
in energy management and identifies common methods and domains. & Need
for uncertainty handling, real-time data integration and hybrid MCDM
with AI, IoT and blockchain.\
\cite{arar_tahir_scientific_2023} & 2023 & 2307 Scopus records &
Bibliometric mapping and science mapping with SciMAT & Optimization of
renewable-energy microgrids & Maps thematic evolution of microgrid,
renewable-energy and optimization research. & Challenges include
stakeholder, technical and financial barriers and renewable-generation
uncertainty.\
\cite{velasquez_intelligence_2023} & 2023 & ref. list: 274 & Bibliometric
mapping of AI/intelligence techniques & Sustainable energy research
across AI application domains & Maps AI-related publication trends,
topics and co-occurrence structures in sustainable energy. & Broad
bibliometric scope; need for more method-specific and decision-oriented
synthesis.\
\end{tabularx}
\end{table*}

Bhosekar and Ierapetritou \cite{Bhosekar2018} review surrogate models for modeling, feasibility
analysis and derivative-free optimization, with emphasis on model
choice, sampling and the use of Kriging and radial-basis-function
models. Díaz-Manríquez et al.Â~\cite{DiazManriquez2016} focus on surrogate-assisted MOO
evolutionary algorithms and classify methods according to how the
surrogate is used inside the evolutionary search algorithm, including
fitness replacement and evolution control. Fernández-Godino \cite{FernandezGodino2023} reviews
multi-fidelity surrogate models and uncertainty quantification, where
low- and high-fidelity models are combined to reduce computational
effort. Starke and da Silva \cite{Starke2025214} discuss
machine-learning-based metamodeling.

In the energy domain, Westermann and Evins \cite{Westermann2019} review surrogate modeling for
sustainable building design and show how surrogates are used for
sensitivity analysis, uncertainty analysis and optimization. Elwy and
Hagishima \cite{Elwy2024} review surrogate-assisted building-performance optimization, including
sampling strategies, design variables, objectives and optimization
workflows, and Zhang et al.Â~\cite{Zhang2026_building} use AI-driven surrogates from data
generation to design selection. Khan et al.Â~\cite{Khan2024_DT} discuss AI-driven
surrogate modeling in HRES from a digital-twin perspective, focusing on
uncertainty and validation. Elsheikh et al.Â~\cite{Elsheikh2019622} review ANN models for
predicting and optimizing solar-energy devices, while Etghani et
al.Â~\cite{Etghani2025}
summarize ANN--Taguchi approaches for prediction and parameter
optimization in engines, thermal systems, solar and wind systems, heat
exchangers and HRES.

Ruan et al.Â~\cite{Ruan2021221} review how machine learning supports power-system
optimization and distinguish prediction, surrogate and hybrid
optimization roles across optimal power flow and economic dispatch.
Khaloie et al.Â~\cite{Khaloie2025} provide a review of learning-based optimal power flow,
including supervised and unsupervised learning, reinforcement learning,
and graph learning. Mohammadi et al.Â~\cite{Mohammadi2024} focus on surrogate models for
optimal power flow, and Tan et al.Â~\cite{Tan2026} review Gaussian processes in power
systems, especially for uncertainty-aware prediction, risk assessment,
optimization and control. Lim et al.Â~\cite{Lim2025} classify deep learning in
power-system decision making into learning to predict, learning to
surrogate and learning to optimize.

Mancò et al.Â~\cite{Manco2024} provide a review of MES optimization and discuss how LP,
MILP, NLP and MINLP formulations, storage, thermal networks, uncertainty
and flexibility affect model complexity. Klemm and Vennemann \cite{Klemm2021} review 145 tools
for MES optimization in mixed-use districts and show that only few tools
satisfy requirements on optimization capability, sector coverage,
spatial scale, temporal resolution and non-financial criteria.
Fattahi et al.Â~\cite{Fattahi2020} review 19 national energy-system models and identify gaps
related to flexibility, electrification, decentralization, and emerging
technologies. Mylonopoulos et al.Â~\cite{Mylonopoulos202332697} review modeling and optimization
of ship energy systems and identify surrogate-assisted optimization as a
promising direction. Agha Kassab et al.Â~\cite{agha_kassab_comprehensive_2024} review microgrid
sizing and energy-management strategies with PV and renewable
integration, while Nallolla et al.Â~\cite{nallolla_multi-objective_2023} focus on MOO
algorithms for HRES microgrids. Salgueiro et
al.Â~\cite{salgueiro_multi-objective_2019}
summarize MOO metaheuristics for microgrid planning and management, and
Malla et al.Â~\cite{malla_sg_optimization_2024} discuss optimization approaches for MES
and Power-to-X conversion. Zhou \cite{Zhou2024__2} reviews renewable-storage
sizing methods and identifies AI-assisted optimization as an emerging
direction. Li et al.Â~\cite{Li2025_hydrogen} review hydrogen-integrated HRES,
simulation tools and AI-driven optimization.
Vahidinasab et al.Â~\cite{vahidinasab_overview_2020} review distribution-network expansion
planning, while Conti et al.Â~\cite{Conti2026} review optimization strategies for
large-scale PV integration in smart distribution networks.
Batista et al.Â~\cite{batista_optimizing_2023} combine bibliometric analysis with a review
of HRES in reverse osmosis desalination.
